\section{Conclusion and discussion}

\subsection{Summary of findings}

We classified the skill level of JIGSAWS Suturing trials as novice, intermediate or expert under the official leave-one-user-out protocol, following the experiment prescribed for this topic and a protocol that fixed every model and comparison before the first result. Simple trial-level kinematic statistics clearly beat the majority class (macro-F1 \FMOne{} against \FMZero{}, exact permutation $p = \PMOne{}$). Video features from sampled frames did not perform above chance on their own, and fusing them with kinematics, by late or early fusion, did not improve on kinematics alone; the pre-declared primary fusion scored \FMThree{}. The focused improvement, motion-texture features in the spirit of the background reference, left macro-F1 essentially unchanged at \FMFour{}. The shortcut audit explains most of this pattern. The transferable signal in the kinematic statistics is largely the pace of the trial, which separates novices from trained surgeons but not intermediates from experts, and with only two experts the model cannot learn what the two have in common, so the expert class is almost never recognised. The whole-frame appearance features identify the surgeon perfectly while carrying no transferable skill signal, and the scene-dynamics features mostly repeat the pace information.

The principal design choices follow from the size of the problem. With eight surgeons, and two in each trained class, we used linear models with strong, nested regularisation and frozen pretrained features rather than trained deep networks, and we reported pooled macro-F1 with user-level bootstrap intervals and an exact user-level permutation test, because per-fold macro-F1 is undefined when each fold contains one class. These choices make the results honest about their uncertainty, at the price of wide intervals that cannot resolve differences smaller than about 0.1 to 0.2 in macro-F1.

\subsection{Strengths and weaknesses}

The main strength of the study is that its conclusions do not depend on choices made after seeing the results. The folds are the official ones, the held-out surgeon never touches any fitted quantity, every model was evaluated once, and negative results are reported alongside positive ones. The audit also shows which information the classifier actually uses, which matters for any clinical use of such a system. Section~\ref{sec:differences} lists four ways in which the approach differs from the background reference and related work; for the conclusions of this study, the most important were the surgeon-level treatment of uncertainty and the direct measurement of the identity and pace shortcuts.

The method also has clear weaknesses. The method works best for what pace can reveal, namely telling a novice from a trained surgeon on a standardised bench-top task, and it is likely to fail whenever trained surgeons differ in quality rather than speed, as the expert class shows. The models are trained on seven surgeons, so a single atypical surgeon can dominate a fold, and the self-reported skill label measures hours of practice rather than observed performance. The video features describe whole frames and do not look at the needle, the thread or the tissue, which is where an observer would see differences in quality that kinematics cannot capture. Finally, equal-weight late fusion is a crude combination rule that lets a weak modality pull down a stronger one.

\subsection{Future directions}

The most effective step would be more surgeons, for example by training on all three JIGSAWS tasks with the held-out surgeon removed from each, or by adding other datasets, because the expert class cannot be learned from one training surgeon. Separating intermediates from experts probably requires features that are less tied to pace, such as gesture-level timing and error events or movement quality normalised by duration, and the identity probe and pace controls introduced here can serve as a check on any such feature. On the video side, features computed in a region around the instrument tips or on the final state of the suture could capture tissue handling and product quality, which neither the kinematics nor our whole-frame descriptors observe. A fusion weight chosen inside the training folds, or a fusion rule that defers to the stronger modality, would also avoid the dilution observed here. Any of these extensions should keep the strict LOUO protocol and the written pre-registration, since the small number of surgeons makes post-hoc choices easy to overfit.

\subsection*{AI tool declaration}

We used Claude Code with two Anthropic models, Claude Fable 5.1 and Claude Opus 5.5, to plan the experiments, write and debug the code, run the experiments, produce the figures and tables, and draft and edit this report and the slides. The group reviewed the code, the results and the text, and is responsible for the content and quality of the submitted work.
